% !TEX TS-program = pdflatex
% !TeX program = pdflatex
% !TEX encoding = UTF-8
% !TeX spellcheck = en_US

\documentclass[11pt, a4paper]{article}
%\usepackage{fullpage}
\usepackage[left=1cm,right=1cm,top=1cm,bottom=2cm]{geometry}
\usepackage[fleqn]{amsmath}
\usepackage{amssymb}
%\usepackage{indentfirst}
\usepackage[T1]{fontenc}
\usepackage[utf8]{inputenc}
\usepackage[french,english]{babel}
\usepackage{txfonts} 
\usepackage[]{graphicx}
\usepackage{multirow}
\usepackage{hyperref}
\usepackage{parskip}
\usepackage{multicol}
\usepackage{wrapfig}
\usepackage{multicol}

\usepackage{turnstile}%Induction symbole

\usepackage{tikz}
\usetikzlibrary{arrows, automata}
\usetikzlibrary{decorations.pathmorphing}

\renewcommand{\baselinestretch}{1}

\setlength{\parindent}{0pt}


\begin{document}
	
	%\pagestyle{empty} 
	\begin{tabular}{ll}
		\multirow{3}{*}{\includegraphics[width=1.5cm]{../img/esi.mlsid.pdf}} 
		& \'Ecole national Supérieure d'Informatique, Algiers, Algeria\\
		& 2CSSID (2024/2025)\\
		& Machine Learning (ML)
	\end{tabular}\\[.25cm]
	\noindent\rule{\textwidth}{1pt}\\[-0.5cm]
	\begin{center}
		{\LARGE \textbf{Workshop 03: Neural Networks' tools (Pytorch)}}
		\begin{flushright}
			Mr. Abdelkrime ARIES
		\end{flushright}
	\end{center}\vspace{-.25cm}
	\noindent\rule{\textwidth}{1pt}
	
	\begin{center}
		\begin{tabular}{|p{.8\textwidth}|}
		\hline
		Get familiar with a tool to create MLPs and train them.  \\
		Learn how to create custom MLPs using this tool. \\
		Understand how this tool works by programming low level MLPs \\
		\hline
	\end{tabular}
	\end{center}

%\section*{Useful links}
%
%\begin{itemize}
%	\item \textbf{Dataset}: {\scriptsize \url{https://archive.ics.uci.edu/dataset/601/ai4i+2020+predictive+maintenance+dataset}} 
%	\item \textbf{Demo}: {\scriptsize \url{https://github.com/projeduc/ESI_ML/blob/main/TP/SD_2022-2023/TP01/TP01_pretraitement_sujet.ipynb}}
%	\item \textbf{MatplotLib}: {\scriptsize \url{https://matplotlib.org/stable/gallery/index.html}}
%	\item \textbf{imbalanced-learn}: {\scriptsize \url{https://imbalanced-learn.org/stable/user_guide.html#user-guide}}
%	\item \textbf{CRISP-DM (web)}: {\scriptsize \url{https://www.ibm.com/docs/en/spss-modeler/18.5.0?topic=dm-crisp-help-overview}}
%	\item  \textbf{CRISP-DM (pdf)}: {\scriptsize \url{https://www.ibm.com/docs/en/SS3RA7_18.5.0/pdf/ModelerCRISPDM.pdf}}
%\end{itemize}


\begin{multicols}{2}
	\section{Data Preparation}
	
	We will use the \textbf{satellite} dataset:
	\begin{itemize} 
		\item Import the training dataset. 
		\item Separate the input data and normalize it by dividing by \textbf{255}. 
		\item Separate the output data and encode it in OneHot format using \textbf{LabelBinarizer} from \textbf{scikit-learn}.
		\item Transform \textbf{X\_train} and \textbf{Y\_train} to \textbf{torch.Tensor}.
		\item Redo the same with test dataset, except do not transform \textbf{Y\_test} to \textbf{torch.Tensor}.
	\end{itemize}
	
	\section{High Level}
	
	In this section, we will create a project to classify satellite images.
	We will try to use high level APIs to create an MLP.
	In PyTorch, the training loop is not handled implicitly; you have to define it yourself.
	But, you can use some other APIs such as \href{https://github.com/Lightning-AI/pytorch-lightning}{pytorch-lightning}.
	
	\subsection{Sequential model}
	
	Create a sequential model (\textbf{nn1}) composed of:
	\begin{itemize} 
		\item a dense layer having a shape (number of features of the past dataset, 10);
		\item a ReLU activation function;
		\item a dense layer having a shape (10, 10);
		\item a ReLU activation function;
		\item a dense layer having a shape (10, number of classes);
		\item and a Softmax activation function with dim=1.
	\end{itemize}
	Print it to get its structure.
	
	\subsection{Model training}
	
	\begin{itemize}  
		\item Use cross entropy as loss function.
		\item Define Adam as an optimization function by passing \textbf{nn1.parameters()} and a learning rate
		\item Train \textbf{nn1} with a loop over a range of iterations of your choice.
		\begin{itemize} 
			\item Apply a forward pass by passing \textbf{X\_train} into \textbf{nn1}.
			\item Calculate the loss.
			\item Set gradients to 0 (\textbf{optimizer.zero\_grad()}).
			\item Apply a backward pass using the instruction \textbf{loss.backward()}.
			\item Apply an optimization step using the instruction \textbf{optimizer.step()}.
			\item Print the loss.
		\end{itemize}
	\end{itemize}
	
	\subsection{Model testing}
	\begin{itemize}  
		\item Generate the estimations of \textbf{X\_test} using \textbf{nn1}.
		\item Transform these estimations into classes using the method \textbf{inverse\_transform} of our \textbf{LabelBinarizer}.
		\item Use these to print a classification report.
	\end{itemize}
	
	\section{High level with a custom class}
	
	We want to create a custom model which adds layers one after the other.
	It verifies if the layers are compatibles.
	Then, when finished, we cannot add any more layers (we lock it).
	
	\subsection{Custom Layer}
	
	\begin{itemize}  
		\item Create a class called \textbf{MyLayer} which inherits from \textbf{torch.nn.Layer}.
		\item Its obligatory arguments are: the number of inputs and the number of outputs.
		\item Its optional arguments are: a boolean for bias (True by default) and a string of activation function to choose from 'relu', 'sigmoid' and 'linear' ('linear' by default).
		\item Use \textbf{assert} to verify that both the number of inputs and outputs are more than 0.
		\item Add the \textbf{forward} method.
		\item You have some unitary tests, adapt them to your model to test it.
	\end{itemize}
	
	\subsection{Custom Net}
	
	\begin{itemize}  
		\item Create a class called \textbf{MyMLP} which inherits from \textbf{torch.nn.Module}.
		\item It has no arguments; but, it contains 2 attributes: a lock initialized to False and a list to store our layers. 
		Do not use \textbf{list}; use \textbf{torch.nn.ModuleList} instead. 
		This is because the latter let the module access the trainable parameters in the list.
		\item Add a method \textbf{add\_layer} which has one argument: a layer. 
		\begin{itemize} 
			\item It raises an exception when the object is locked.
			\item It raises an exception when the previous layer's output is not as the current layer's input.
			\item In case this is the first layer to add, no need to check.
			\item It adds the layer when all is right
			\item Then, it returns \textbf{self}; this allows us to call this method in one line 
		\end{itemize}
		\item Add a method \textbf{compile} which has these optional arguments: number of inputs of the output layer (=1), number of outputs of the output layer (=1), bias of the output layer (=True), a boolean indicating if it is a multi-class classification (=False), and learning rate (=1.).
		\begin{itemize} 
			\item If there are past layers, it modifies the number of inputs of the output layer to be the output of the last layer.
			\item The output layer will have a sigmoid activation function
			\item If it is multi-class and the number of outputs is more than 1, this function must be softmax.
			\item The loss function must be BCE.
			\item If it is multi-class and the number of outputs is more than 1, The loss function must be CE.
			\item The optimizer must be Adam, taking this model's parameters and the bias.
			\item Lock the model.
		\end{itemize}
		\item Add the \textbf{forward} method taking \textbf{X} as argument.
		\item Add the \textbf{backward} method taking \textbf{X, Y} as arguments:
		\begin{itemize} 
			\item It applies a forward pass.
			\item It calculates a loss using the output and \textbf{Y}.
			\item It initializes the gradients of the optimizer to zero (\textbf{.zero\_grad()}).
			\item It applies a backward pass on the loss.
			\item It applies a step of the optimizer.
			\item It returns the loss as numpy (use \textbf{.detach().numpy()}) .
		\end{itemize}
		\item Add a \textbf{fit} method taking \textbf{X, Y, epochs} as arguments.
		For each iteration, it shows the cost.
		\item Override \textbf{\_\_call\_\_} which returns the backward pass's result.
	\end{itemize}
	
	\subsection{Model training} 
	
	\begin{itemize}  
		\item Create a model (\textbf{nn2}) composed of:
		\begin{itemize} 
			\item a hidden layer having a shape (number of features of the past dataset, 10) and ReLU activation function;
			\item a hidden layer having a shape (10, 10) and a ReLU activation function;
			\item an output layer having a shape (10, number of classes) for multi-class classification.
		\end{itemize}
		\item  Print it to get its structure.
		\item Fit it on \textbf{satellite} train dataset.
	\end{itemize}
	
	\subsection{Model testing} 
	
	\begin{itemize}  
		\item Generate the estimations of \textbf{X\_test} using \textbf{nn2}.
		\item Transform these estimations into classes using the method \textbf{inverse\_transform} of our \textbf{LabelBinarizer}.
		\item Use these to print a classification report.
	\end{itemize}
	
	\section{Low Level}
	
	\subsection{Activation functions}
	
	\begin{itemize}  
		\item Create classes \textbf{SimpleSigmoid}, \textbf{SimpleReLU} and \textbf{SimpleSoftmax} which inherit from \textbf{torch.nn.Module}.
		\item No constructor is required. 
		\item Add the \textbf{forward} method which takes \textbf{X} as argument.
	\end{itemize}
	
	\subsection{Loss functions}
	
	\begin{itemize}  
		\item Create classes \textbf{SimpleBCE} and \textbf{SimpleCE} which inherit from \textbf{torch.nn.Module}.
		\item No constructor is required. 
		\item Add the \textbf{forward} method which takes \textbf{X, Y} as arguments.
	\end{itemize}
	
	\subsection{Optimization functions}
	
	\begin{itemize}  
		\item Create a class \textbf{SimpleGD} which inherits from \textbf{optim.Optimizer}.
		\item The constructor takes \textbf{params} and \textbf{lr} as arguments. 
		\item It passes them to the super class as \textbf{super().\_\_init\_\_(params, defaults=\{'lr': lr\})}
		\item Add the \textbf{step} method which takes no arguments. 
		This function must update all parameters based on their gradients and their learning rate.
		\begin{itemize}  
			\item \textbf{self.groups} contains a list of parameters' groups.
			\item each group is a dictionary containing, among others, a key '\textbf{lr}' and another '\textbf{params}'.
			\item the value of '\textbf{lr}' is the learning rate of this group of parameters.
			\item the value of '\textbf{params}' is a list of tensors. 
			Since these tensors are 'trainable', they have an attribute \textbf{data} containing theirs current values, and another \textbf{grad} containing their current gradients.
			{\color{red} These gradients are coming from \textbf{loss.backward()} which will calculate the gradients of all the trainable parameters contributing to the loss and for each stores the gradient in its \textbf{grad} attribute}.
		\end{itemize}
	\end{itemize}
	
	\subsection{Custom Layer}
	
	\begin{itemize}  
		\item Copy \textbf{MyLayer} and rename it \textbf{SimpleLayer}, but it must inherit from \textbf{torch.nn.Module}.
		\item Create an attribute \textbf{W} of type \textbf{nn.parameter.Parameter} initialized to zeros.
		\item Create an attribute \textbf{W} initialized to zeros. 
		It can be only transformed to \textbf{nn.parameter.Parameter} if bias=True.
		In this case, it will be subscribed into the module's parameters; thus can be recovered using the method \textbf{parameters}.
		\item Replace all activation functions by our simple versions.
		\item Add a method \textbf{randomize} which randomizes all trainable parameters using a normal law (\textbf{torch.normal}) having a mean of 0 and a standard deviation of 1.
		Try to update their \textbf{data} attribute.
	\end{itemize}
	
	\subsection{Custom Net}
	
	\begin{itemize}  
		\item Copy \textbf{MyMLP} and rename it into \textbf{SimpleMLP}.
		\item Modify it to handle our new structures.
		\item Add a method \textbf{randomize} which randomizes all layers. 
	\end{itemize}
	
	\subsection{Model training} 
	\begin{itemize}  
		\item Create a model (\textbf{nn3}) composed of:
		\begin{itemize} 
			\item a hidden layer having a shape (number of features of the past dataset, 10) and ReLU activation function;
			\item a hidden layer having a shape (10, 10) and a ReLU activation function;
			\item an output layer having a shape (10, number of classes) for multi-class classification.
		\end{itemize}
		\item Randomize it.
		\item Print its parameters. The method \textbf{parameters} returns an iterator, so you have to cast it into a list.
		\item Fit it on \textbf{satellite} train dataset.
	\end{itemize}
	
	\subsection{Model testing} 
	
	\begin{itemize}  
		\item Generate the estimations of \textbf{X\_test} using \textbf{nn3}.
		\item Transform these estimations into classes using the method \textbf{inverse\_transform} of our \textbf{LabelBinarizer}.
		\item Use these to print a classification report.
	\end{itemize}
	
	
\end{multicols}

\section*{Appendix}

\begin{itemize}  
	\item pandas.read\_csv: \url{https://pandas.pydata.org/docs/reference/api/pandas.read_csv.html}
	\item sklearn.preprocessing.LabelBinarizer: \url{https://scikit-learn.org/1.5/modules/generated/sklearn.preprocessing.LabelBinarizer.html}
	\item torch.Tensor: \url{https://pytorch.org/docs/main/tensors.html}
	\item torch.nn.Sequential: \url{https://pytorch.org/docs/stable/generated/torch.nn.Sequential.html}
	\item torch.nn.Linear: \url{https://pytorch.org/docs/stable/generated/torch.nn.Linear.html#torch.nn.Linear}
	\item torch.nn.ReLU: \url{https://pytorch.org/docs/stable/generated/torch.nn.ReLU.html#relu}
	\item torch.nn.Softmax: \url{https://pytorch.org/docs/stable/generated/torch.nn.Softmax.html#softmax}
	\item torch.optim.Adam: \url{https://pytorch.org/docs/stable/generated/torch.optim.Adam.html#adam}
	\item torch.nn.CrossEntropyLoss: \url{https://pytorch.org/docs/stable/generated/torch.nn.CrossEntropyLoss.html#crossentropyloss}
	\item sklearn.metrics.classification\_report: \url{https://scikit-learn.org/1.5/modules/generated/sklearn.metrics.classification_report.html}
	\item assert: \url{https://docs.python.org/3/reference/simple_stmts.html#assert}
	\item torch.nn.Sigmoid: \url{https://pytorch.org/docs/stable/generated/torch.nn.Sigmoid.html#torch.nn.Sigmoid}
	\item Exception: \url{https://docs.python.org/3/library/exceptions.html#Exception}
	\item torch.nn.Module: \url{https://pytorch.org/docs/stable/generated/torch.nn.Module.html#torch.nn.Module}
	\item torch.nn.ModuleList: \url{https://pytorch.org/docs/stable/generated/torch.nn.ModuleList.html#torch.nn.ModuleList}
	\item torch.nn.BCELoss: \url{https://pytorch.org/docs/stable/generated/torch.nn.BCELoss.html#bceloss}
	\item torch.exp: \url{https://pytorch.org/docs/stable/generated/torch.exp.html#torch-exp}
	\item torch.where: \url{https://pytorch.org/docs/stable/generated/torch.where.html#torch-where}
	\item torch.Tensor.view: \url{https://pytorch.org/docs/stable/generated/torch.Tensor.view.html}
	\item torch.sum: \url{https://pytorch.org/docs/stable/generated/torch.sum.html#torch.sum}
	\item torch.mean: \url{https://pytorch.org/docs/stable/generated/torch.mean.html#torch-mean}
	\item torch.log: \url{https://pytorch.org/docs/stable/generated/torch.log.html#torch-log}
	\item torch.optim.Optimizer: \url{https://pytorch.org/docs/stable/optim.html#torch.optim.Optimizer}
	\item torch.nn.parameter.Parameter: \url{https://pytorch.org/docs/stable/generated/torch.nn.parameter.Parameter.html#torch.nn.parameter.Parameter}
	\item torch.zeros: \url{https://pytorch.org/docs/stable/generated/torch.zeros.html#torch-zeros}
	\item torch.normal: \url{https://pytorch.org/docs/stable/generated/torch.normal.html#torch-normal}
	\item torch.matmul: \url{https://pytorch.org/docs/stable/generated/torch.matmul.html#torch-matmul}
\end{itemize}


\end{document}
